\noindent\fcolorbox{accent!40}{tint}{\begin{minipage}{0.965\linewidth}\vspace{0.3em}
\textbf{\sffamily 本部分的内容与读法}\quad 每个算例一节，每节先给出“读图”（逐行说明每一行画的是什么、颜色代表什么），再给出“定量读数”（从图中各场直接算出的数，脚本 \texttt{scripts/viz\_stats.py}，结果 \texttt{results/res\_viz\_stats.json}），最后是“物理解释”与“伪影与局限”。图中的场全部只用示踪点位置（加初始密度）计算，标注为 DNS 的行仅作对照。通用约定如下：
\begin{itemize}[leftmargin=1.4em,itemsep=0.1em]
\item \textbf{物质标签图}：按初始归属给示踪点着色（红、蓝两类物质），黑点为初始位于某条物质线或界面上的示踪点。它显示的是\emph{平流与物质输运}本身：红色物质跑到哪里，就有多少物质被输运过去。
\item \textbf{扩散标量图}：初始为一条光滑锋面的被动标量，在网络上按质量形式的拉格朗日扩散方程演化（T8）。锋面被拉长、变薄并展宽，显示的是\emph{平流与扩散的共同作用}，即混合。
\item \textbf{拉伸场 $\ln\lambda$}：链式张量估计的初始标签等值线族的切向拉伸，红为拉长、蓝为缩短；\textbf{压缩场 $-\ln J$}：账本给出的面积压缩，橙为压缩、紫为膨胀。二者之和就是界面上标量梯度的增长（Nanson 恒等式），即\emph{掺混的两个来源}。
\item \textbf{各向异性 $A=\sigma_{\max}(F)/\lambda$}：大值表示标签等值线不在最大拉伸方向上，链式张量在这里对误差敏感。
\end{itemize}
\vspace{0.2em}\end{minipage}}

\section{双环流（不可压，解析速度场）}
\paragraph{读图。}第一行：左环流的物质为红、右环流为蓝，黑点是初始物质线 $x_0=0.5$；第二行：初始锋面在 $x=1$ 的扩散标量（$\kappa=10^{-3}$）；第三行：只凭位置的步涡量 $\hat\omega=2A\,\mathrm{Im}\,p/(\bar A\Delta t)$；第四行：$t=5$ 时由位置差分得到的速度矢量，以及 $t=10$ 时的拉伸场、各向异性与压缩场。四列为 $t=0,3.4,6.6,10$。
\paragraph{定量读数。}左环流的物质进入右半区的比例依次为 0、14.6\%、20.0\%、32.9\%：分界线 $x\approx1$ 的周期摆动每次把一个“瓣”的物质卷入对侧。扩散标量处于混合区（$0.1<c<0.9$）的示踪点比例为 11.3\%、18.5\%、35.1\%、59.2\%，按质量加权的方差比为 1、0.93、0.77、0.56。步涡量在 $t\approx6.6$ 的 1\%–99\% 分位为 $-2.5$ 到 $1.7$，两个环流的符号相反。$t=10$ 时拉伸场的中位数为 0.96，5\% 与 95\% 分位为 $-0.74$ 与 3.14，22.6\% 的示踪点 $|\ln\lambda|>2$；跟踪的物质线上平均拉伸为 1.708（逐点精确值 1.711）。压缩场的 95\% 分位只有 0.06。
\paragraph{物理解释。}这是典型的瓣输运与 Batchelor 型混合：物质每个周期被折叠一次，物质线被拉长、折叠成 S 形，扩散标量的锋面沿物质丝被拉薄，梯度变陡，扩散随之加速，所以方差在后半段衰减得更快（$t=6.6\to10$ 的方差比下降 0.21，前 3.4 个时间单位只下降 0.07）。各向异性的中位数为 1.36、95\% 分位为 14.4，且 $|\ln\lambda|$ 与 $A$ 的秩相关为 $-0.45$：拉伸最强的地方，等值线恰好沿着最大拉伸方向（$A\approx1$）；$A$ 大的地方是等值线被侧向剪切、切向反而不长的区域。
\paragraph{伪影与局限。}不可压流的压缩场应恒为零，图中 $\pm0.06$ 的斑点是账本在强折叠物质丝上的离散误差，与 T2 的 95 分位（5.7\%）一致。
\begin{figure}[p]\centering\includegraphics[width=\linewidth]{viz_dg.pdf}
\caption{双环流的场图。第一行：物质标签与物质线 $x_0=0.5$；第二行：网络扩散标量（质量形式）；第三行：只凭位置的步涡量；第四行：位置差分速度、拉伸场、各向异性与压缩场。}\end{figure}

\clearpage\section{Lamb–Oseen 共转涡对（不可压）}
\paragraph{读图。}第一行：$x<0$ 的物质为红，黑点是绕左涡、半径 1.5 的物质圆环；第二行：初始锋面在 $x=0$ 的扩散标量（$\kappa=0.02$）；第三行：步涡量；第四行：速度矢量与末时刻（$t=6\pi$）以绕左涡圆周为切向的拉伸场、各向异性和压缩场。四列为 $t=0,6.1,12.7,18.8$。
\paragraph{定量读数。}原在左半区的物质进入右半区的比例为 0、15.9\%、31.1\%、40.5\%，这主要是涡对整体公转带来的输运，而不是混合；混合区比例为 6.5\%、14.8\%、22.2\%、31.7\%，方差比为 1、0.93、0.85、0.76，混合明显慢于双环流。步涡量的峰值为 1.81（解析峰值 $\Gamma/\pi a^2=2$，三角形平均使峰值略低）。物质圆环到末时刻仍近圆：相对其中心的半径的相对标准差为 4.6\%；圆环上的局部拉伸在 $-2.31$ 到 $0.55$ 之间（5\%–95\% 分位），平均 $-0.59$。圆环上各向异性的中位数为 14.1，全场只有 2.0。
\paragraph{物理解释。}涡对的输运以整体旋转为主，两臂螺旋把物质卷绕进去，但涡核本身保持相干，所以扩散标量的锋面只是被卷绕、没有被大量拉薄，方差衰减慢。物质圆环在差动旋转下总长几乎不变，但局部被另一个涡的应变时而拉长、时而缩短，这就是第一部分“圆环拉伸约为 0”被推翻的直观原因：平均为负，局部在 $-2.3$ 到 $0.6$ 之间。
\paragraph{伪影与局限。}圆环恰处在强剪切中（$A\approx14$），链式张量的误差在这里被放大（第一部分：$N=8000$ 时界面平均 $-0.55$，精确 $-0.31$，随分辨率二阶收敛）；压缩场 95\% 分位 0.13 也来自涡周星的非仿射变形，自适应锚定把这一区域的密度 95 分位从 13\% 降到 7\%。
\begin{figure}[p]\centering\includegraphics[width=\linewidth]{viz_lo.pdf}
\caption{Lamb–Oseen 共转涡对的场图。各行同双环流；第四行的拉伸场以绕左涡的圆周为切向。}\end{figure}

\clearpage\section{Sod 激波管（可压，精确轨迹）}
\paragraph{读图。}第一行：驱动气体（$x_0<0$，红）与被驱动气体（蓝），黑点为物质线 $x_0=0.1$；第二行：只凭位置的账本密度；第三行：账本密度的相对误差（色标 $\pm5\%$）；第四行：初始锋面在接触面的扩散标量（$\kappa=10^{-3}$）；第五行：C 通道检出的穿越时刻（画在初始位置上）、$t=0.14$ 的步膨胀率 $\hat\theta$、同一时刻的步涡量（理论上应为零）与速度矢量。四列为 $t=0.04,0.105,0.175,0.24$。
\paragraph{定量读数。}由检出的穿越时刻对初始位置线性拟合，得激波速度 1.755（精确 1.752）；被检出穿越的示踪点占 35\%。账本给出的波后平台密度为 0.26557，接触面左侧为 0.42632，与精确解的 0.26557、0.42632 一致到 5 位；相对误差超过 1\% 的示踪点只占 5.1\%，全部在激波与接触间断两侧的模板内。扩散锋面的 10\%–90\% 宽度依次为 0.061、0.071、0.077、0.080。激波处步膨胀率达到 $-151$，同一步的跨激波三角形上伪涡量的最大幅值为 23，约为膨胀率的 15\%。
\paragraph{物理解释。}物质标签图直接显示接触面就是两种气体的物质边界，以速度 0.93 右移；物质线 $x_0=0.1$ 在激波扫过时被一次性压缩（压缩 0.754），之后随流平移，拉伸恒为零——平面激波对掺混的贡献完全来自压缩。扩散锋面的宽度增长随时间变慢，是一维扩散 $\sqrt{\kappa t}$ 的表现；锋面在接触面两侧的密度不同（0.43 与 0.27），质量形式正确计入了这一密度差（体积形式误差大一个数量级，见第二部分）。
\paragraph{伪影与局限。}第五行第三个子图里沿激波的条带是伪涡量：跨激波三角形的三点映射不是单轴的，$\mathrm{Im}\,p\neq0$。这是“跨激波只报告积分量、用合并块计算转角”的直观理由；合并块后伪转角的 95 分位从 0.092 rad 降到 0.0057 rad。
\begin{figure}[p]\centering\includegraphics[width=\linewidth]{viz_sod.pdf}
\caption{Sod 激波管的场图。第一至四行：物质标签、账本密度、密度误差、扩散标量；第五行：穿越时刻、步膨胀率、伪涡量与速度。}\end{figure}

\clearpage\section{柱面爆炸波（可压，精确轨迹）}
\paragraph{读图。}第一行：$r_0\le0.2$ 的高密度内部气体为红，黑点为物质圆环 $r_0=0.5$；第二行：账本密度（色标 0–6）；第三行：初始锋面在 $r=0.5$ 的扩散标量（$\kappa=2\times10^{-4}$）；第四行：穿越时刻、$t=25$ 的拉伸场（切向取圆周方向）与压缩场、$t=13$ 的速度矢量。四列为 $t=1,9,17,25$。
\paragraph{定量读数。}对检出的穿越时刻与初始半径做双对数拟合，指数为 0.54（理论的 $r_s\propto t^{1/2}$ 给出 0.5）。末时刻外部气体被压进 $r\approx0.94$–$1.00$ 的薄壳，账本密度的 99.5\% 分位为 5.73（壳内精确峰值约 6）；壳内压缩的中位数 1.48。内部气体强烈膨胀：中心区 $-\ln J$ 为 $-2.68$，即密度降到初始的 $1/14.6$，内部气体的中位半径从 0.2 增大到 0.54。圆环上拉伸 0.576（精确 0.578）、压缩 0.637（精确 0.638）。扩散标量的混合区比例为 6.1\%、21.0\%、43.9\%、55.5\%，方差比为 1、0.85、0.65、0.56。
\paragraph{物理解释。}扩散标量的环形锋面先被激波压进薄壳而变陡（$t=9$ 混合区已达 21\%），随后在膨胀流中被拉长并展宽；方差在激波扫过锋面的阶段衰减最快。圆环同时受到两种作用：激波一次性的压缩（使梯度变陡）与之后随流膨胀的径向拉伸（使圆周变长），所以界面梯度的增长是“先压缩、后被膨胀部分抵消”。拉伸场在全场为径向自相似的均匀拉伸，没有局部结构，与流动的一维性质一致。
\paragraph{伪影与局限。}薄壳只有约三个示踪点间距厚，激波正扫过的示踪点的星跨在激波上，这是账本 95 分位（32\%）的来源；第四部分的断轨检验中，爆炸波的薄壳也是唯一仍有大误差的区域。
\begin{figure}[p]\centering\includegraphics[width=\linewidth]{viz_blast.pdf}
\caption{柱面爆炸波的场图。第一至三行：物质标签与圆环、账本密度、扩散标量；第四行：穿越时刻、拉伸场、压缩场与速度。}\end{figure}

\clearpage\section{激波–等熵涡（可压，DNS）}
\paragraph{读图。}第一行：初始涡核（$r_0<R_v$）的物质为红，黑点为 $r_0=2R_v$ 的物质圆环；第二行：账本密度；第三行：账本密度相对 DNS 的误差（色标 $\pm5\%$）；第四行：只凭位置的步涡量；第五行：DNS 涡量（对照）。四列为 $t=0,0.2,0.4,0.6$。
\paragraph{定量读数。}涡核物质的质心位移为 0、0.0004、0.060、0.149：激波在 $t\approx0.2$ 到达涡核之前涡核几乎不动，之后被波后流动带走。圆环的平均压缩为 0、0.12、0.42、0.40，激波扫过后基本不变。账本密度相对 DNS 的中位误差为 0.09\%、0.22\%、0.32\%，95 分位为 1.1\%、2.4\%、3.2\%。$t=0.4$ 时误差最大的 5\% 中有 27\% 位于近期的激波模板内，末时刻有 93\% 位于涡核内。步涡量与 DNS 涡量在示踪点上的相关系数为 0.994、0.998、0.999（三个时刻）。
\paragraph{物理解释。}激波扫过等熵涡时，涡核物质保持相干、被压缩成椭圆并随波后流动输运；涡量几乎不被激波改变（等熵涡没有密度梯度与压力梯度的交叉，斜压沉积很小），所以位置涡量与 DNS 涡量在激波前后都高度一致。误差的位置随时间迁移：激波经过时集中在激波带，之后集中在涡核——后者是涡核内锚定星的非仿射变形，第三部分的 DNS 加密已确认它与 DNS 网格无关。
\paragraph{伪影与局限。}激波正扫过涡核的帧上，第四行可以看到沿激波的伪涡量条带，第五行的 DNS 涡量里没有；在这些三角形上应只报告积分量。
\begin{figure}[p]\centering\includegraphics[width=\linewidth]{viz_sv.pdf}
\caption{激波–等熵涡的场图。第一至三行：涡核物质与圆环、账本密度、密度误差；第四、五行：只凭位置的步涡量与 DNS 涡量。}\end{figure}

\clearpage\section{激波–重气泡（可压，无粘 DNS）}
\paragraph{读图。}第一行：重气体示踪点（红）与界面示踪点（黑）；第二、三行：账本密度与 DNS 密度；第四、五行：网络扩散标量（质量形式）与带同一扩散系数的 DNS 标量；第六行：$t=1.4$ 的拉伸场（切向取初始圆周方向）、压缩场、各向异性与 $t=1.38$ 的步涡量。四列为 $t=0,0.6,1.0,1.4$。
\paragraph{定量读数。}重气体的流向范围从 $[0.52,0.98]$ 变为 $[0.71,0.98]$（$t=0.6$，被压扁）、$[0.87,1.15]$、$[1.07,1.42]$（$t=1.4$，卷成肾形并整体下游输运）；最高点从 0.24 先降到 0.22，再被卷起抬到 0.30。账本密度相对 DNS 的中位误差为 0.28\%、0.58\%、0.67\%，末时刻误差最大的 5\% 中有 87\% 位于距界面 0.03 以内。网络扩散标量与 DNS 标量的相对 $L^2$ 差为 1.4\%、2.3\%、7.3\%；混合区比例几乎相同（网络 23.7\%、33.1\%、40.2\%，DNS 24.0\%、32.9\%、40.4\%）。$t=1.4$ 时界面上的拉伸在 $-1.50$ 到 1.53 之间（5\%–95\% 分位），平均 0.29；压缩在界面上平均 0.47，在重气体内部 0.63，在轻气体中 0.58。界面上各向异性的中位数为 5.0，远离界面为 2.1。步涡量的 1\%–99\% 分位为 $-18.2$ 到 3.3，以负值（上半平面顺时针）为主。
\paragraph{物理解释。}物质标签图给出了输运的完整图像：激波先把气泡沿流向压扁（$t=0.6$），斜压沉积的涡片随后把重气体卷成涡对（$t=1.0$–1.4），重气体沿螺旋臂进入涡中。拉伸场与压缩场把掺混分成两部分：压缩是激波给的、几乎均匀（重气体内 0.63）；拉伸集中在螺旋臂上，有正有负——卷起的臂被拉长（红），与激波大致平行的界面段被切向缩短（蓝）。这正是第三部分“激波扫过后界面先被切向缩短、抵消大部分压缩”的空间图像。步涡量以负值为主，对应上半气泡顺时针的斜压环量（第三部分：位置形式的上半气泡环量 0.266，与 DNS 相差 1.5\%，与 Samtaney–Zabusky 模型相差 8\%）。网络扩散标量与 DNS 标量的混合区几乎一致，说明质量形式的网络扩散在卷起的螺旋中仍然可靠。
\paragraph{伪影与局限。}界面附近的密度误差主要来自 DNS 本身的数值扩散（第二部分的 DNS 加密）；卷起后的涡片上各向异性大（$A\approx5$），链式张量在这里误差被放大，且无粘 DNS 的界面带不收敛（信息极限）。
\begin{figure}[p]\centering\includegraphics[width=\linewidth]{viz_sb.pdf}
\caption{激波–重气泡（无粘 DNS）的场图。第一至五行：重气体与界面、账本密度、DNS 密度、网络扩散标量、DNS 扩散标量；第六行：拉伸场、压缩场、各向异性与步涡量。}\end{figure}

\clearpage\section{粘性激波–重气泡（$Re\approx1300$）}
\paragraph{读图。}第一行：重气体与界面，$t=0.6,1.0,1.4$；第二行：账本密度；第三行：$t=1.4$ 的密度相对误差（对粘性 DNS，色标 $\pm10\%$）、拉伸场与压缩场。
\paragraph{定量读数。}重气体的整体输运与无粘算例相近（$t=1.4$ 流向范围 $[1.08,1.40]$，最高点 0.297），但界面光滑、没有细小卷起。账本密度相对粘性 DNS 的中位误差为 0.27\%、0.43\%、0.48\%，误差最大的 5\% 中有 79\% 位于界面附近。界面上的拉伸在 $-1.12$ 到 1.30 之间，平均 $-0.02$；压缩在界面上平均 0.32，在重、轻气体内部为 0.63 与 0.58；界面各向异性的中位数为 3.7（无粘为 5.0）。
\paragraph{物理解释。}物理粘性抑制了涡片的卷起，界面整体没有被拉长（平均拉伸接近零），所以到 $t=1.4$ 界面上梯度的增长几乎全部来自压缩。与无粘算例对比，两者的整体输运相同，差别在界面的小尺度结构：这是“压缩与拉伸的分配依赖雷诺数”的直观证据。界面压缩（0.32）低于两侧内部（约 0.6），是因为物理扩散在界面处把两种气体的密度混合，界面上的物质元在压缩后又随混合膨胀。
\begin{figure}[p]\centering\includegraphics[width=\linewidth]{viz_sbv.pdf}
\caption{粘性激波–重气泡（$\Delta x=0.0035$）的场图。第一、二行：重气体与界面、账本密度；第三行：密度误差、拉伸场与压缩场。}\end{figure}

\clearpage\section{单模 Richtmyer–Meshkov（可压，DNS）}
\paragraph{读图。}第一行：重气体（$x>0.5+a_0\cos2\pi y$，红）与界面示踪点；第二行：账本密度；第三行：步涡量；第四行：$t=2$ 的密度误差、拉伸场（切向取初始界面方向）、压缩场与 $t=0.4$ 的速度矢量。四列为 $t=0,0.6,1.2,2.0$。
\paragraph{定量读数。}激波扫过后界面以 $\Delta u\approx0.51$（速度矢量的中位数）整体右移，扰动振幅先被压缩、随后线性增长（第二部分：增长率 0.042，Richtmyer 模型 0.043）。$t\approx1.2$ 时距界面 0.03 以内的示踪点承载了 65\% 的 $|\hat\omega|$ 积分：涡量集中在界面上的斜压涡片中。$t=2$ 时界面压缩平均 0.80，重、轻气体内部为 0.76 与 0.78；拉伸在尖钉处（$y\approx0$）为 $+0.34$，气泡顶（$y\approx0.5$）为 $+0.07$，侧翼（$y\approx0.25$）为 $-0.12$。远离界面的密度中位误差为 0.44\%。
\paragraph{物理解释。}重气体的尖钉在 $y=0$ 处伸入轻气体、轻气体的气泡在 $y=0.5$ 处伸入重气体，界面上的涡片驱动这一增长。拉伸场显示界面被拉长的位置是尖钉与气泡顶，而侧翼被缩短；但到 $ka\approx0.6$ 拉伸的幅度（0.1–0.3）仍远小于激波的压缩（0.8），所以这一阶段界面梯度的增长约 94\% 来自压缩。压缩场在两侧气体中几乎相同，说明激波压缩是整体性的。
\paragraph{伪影与局限。}紧贴界面（$|\phi_0|<0.02$）的密度“误差”中位数高达 35\%，但这不是账本的误差：DNS 在接触间断上有数值扩散，把两种气体的密度混在一起（与激波–重气泡的 DNS 加密结论相同），界面上的 DNS 密度不能作参照；远离界面账本与 DNS 一致到 0.44\%。
\begin{figure}[p]\centering\includegraphics[width=\linewidth]{viz_rm.pdf}
\caption{单模 Richtmyer–Meshkov 的场图。第一至三行：重气体与界面、账本密度、步涡量；第四行：密度误差、拉伸场、压缩场与速度。}\end{figure}

\clearpage\section{小结：图中能看到的五个机理}
\begin{enumerate}
\item \textbf{平流与物质输运}：物质标签图把输运直接画了出来（瓣输运、涡对卷绕、激波后整体平移、气泡卷成涡对、RM 的尖钉与气泡），并可按质量计量（双环流 32.9\%、Lamb–Oseen 40.5\% 的物质跨过中线）。
\item \textbf{压缩}：激波在物质上留下一次性的、几乎均匀的压缩（Sod 0.754、激波–重气泡约 0.6、RM 约 0.8），膨胀则表现为负的 $-\ln J$（爆炸波中心 $-2.7$）。
\item \textbf{拉伸}：只在有涡或剪切的地方出现，且有正有负——双环流的物质丝、Lamb–Oseen 的涡周、激波–重气泡的螺旋臂、RM 的尖钉。
\item \textbf{扩散与混合}：扩散标量在平流拉薄的地方混合最快（双环流后半段、爆炸波激波扫过锋面时），网络质量形式与 DNS 标量的混合区几乎一致。
\item \textbf{方法的边界}：跨激波三角形的伪涡量条带、涡周与涡片上的高各向异性、接触间断上 DNS 的数值扩散，都在图中直接可见。
\end{enumerate}
